\documentclass[11pt,letterpaper]{article}

\usepackage[utf8]{inputenc}
\usepackage[T1]{fontenc}
\usepackage{helvet}
\renewcommand{\familydefault}{\sfdefault}
\usepackage{microtype}
\usepackage{xcolor}
\usepackage[letterpaper, margin=0.5in, top=0.4in, bottom=0.35in]{geometry}
\usepackage{enumitem}
\usepackage{tabularx}
\usepackage{titlesec}
\usepackage[hidelinks]{hyperref}
\hypersetup{colorlinks=true, urlcolor=blue}

\pagestyle{empty}
\setlength{\parindent}{0pt}
\setlength{\parskip}{0pt}

\titleformat{\section}{\normalsize\bfseries\MakeUppercase}{}{0em}{}[\vspace{0.5pt}\hrule\vspace{1.5pt}]
\titlespacing*{\section}{0pt}{4pt}{1.5pt}

\setlist[itemize]{nosep, leftmargin=12pt, itemsep=0pt, topsep=0.5pt}

\begin{document}

\begin{center}
{\large\bfseries JORGE J.\ ORTIZ, Ph.D.}\\[1pt]
{\small Associate Professor (with Tenure) \textbullet\ Electrical and Computer Engineering \textbullet\ Rutgers University}\\[1pt]
{\small \href{mailto:jorge.ortiz@rutgers.edu}{jorge.ortiz@rutgers.edu}
\quad\textbar\quad \href{http://jorgeortizphd.info}{jorgeortizphd.info}
\quad\textbar\quad 617-784-6550}
\end{center}

\vspace{1pt}

\section{Research Focus}

{\small Multimodal machine learning for sensing and human-state inference, with emphasis on interpretable fusion of heterogeneous sensor streams (physiological, inertial, audio), time-series modeling, and human-centered AI for health and wearable systems.
Director of the Sensing and Reasoning (SnR) Lab; 40+ peer-reviewed publications; 12+ issued U.S.\ patents.}

\section{Education}

{\small\begin{tabularx}{\textwidth}{@{}X r@{}}
\textbf{Ph.D., Computer Science} -- University of California, Berkeley & 2013 \\
\textbf{M.S., Computer Science} -- University of California, Berkeley & 2010 \\
\textbf{B.S., Computer Science} -- Massachusetts Institute of Technology & 2003 \\
\end{tabularx}}

\section{Positions}

{\small\begin{tabularx}{\textwidth}{@{}X r@{}}
\textbf{Associate Professor (with Tenure)}, ECE, Rutgers University & 2025 -- Present \\
\textbf{Director, Sensing and Reasoning Lab}, Rutgers University & 2018 -- Present \\
\textbf{Leadership Research Faculty Fellow}, Rutgers Office for Research & 2026 \\
\textbf{AI \& Computer Vision Lead}, New York Yankees (Baseball Operations) & 2019 -- Present \\
\textbf{Research Staff Member}, IBM Research & 2013 -- 2018 \\
\end{tabularx}}

\section{Expertise Relevant to This Proposal}

{\small\begin{itemize}
\item \textbf{Interpretable multimodal sensor fusion} -- Developed GeXSe, an interpretable deep generative framework that fuses heterogeneous non-invasive sensor streams with human-understandable modality-relevance explanations via multi-branch MLP and Fast Fourier Convolution (\textit{IEEE Sensors J.}, 2025).
\item \textbf{Robust temporal activity recognition} -- Developed VCHAR, a framework for complex activity recognition under imprecise temporal supervision that exposes sensor contributions through interpretable cross-modal explanations (\textit{arXiv} 2407.03291, 2024).
\item \textbf{Health \& wearable sensing} -- Co-PI on NIH CAMERA platform (\$5.4M) integrating neural, physiological, behavioral, and environmental sensors for continuous mental-health state prediction.
\item \textbf{IoT device modeling} -- Designed ML pipelines for identifying unknown IoT devices from network behavior (\textit{ACM IoTDI}, 2019; Best Paper Finalist).
\end{itemize}}

\section{Selected Grants (\$32.4M Total Portfolio)}

{\small\begin{itemize}
\item \textbf{NIH CAMERA Platform}: Multimodal ecological research for mental health, Co-PI (\$5.4M), 2024--2029
\item \textbf{NSF ERC -- Center for Smart Streetscapes}: Rutgers Site PI (\$26M total), 2022--present
\item \textbf{NSF IUCRC CRAIG}: Responsible AI and Governance, PI at Rutgers (\$1.5M), 2027--2029
\item \textbf{NSF ReDDDoT Phase 2}: Urban AI in Harlem, Co-PI (\$1.45M), 2024--present
\end{itemize}}

\section{Selected Publications}

{\small\begin{itemize}
\item Y.\ Sun, N.\ Salami Pargoo, \textbf{J.\ Ortiz}, ``GeXSe: Generative Explanatory Sensor System -- An Interpretable Deep Generative Model for Human Activity Recognition,'' \textit{IEEE Sensors J.}, vol.~25, no.~7, 2025.

\item Y.\ Sun et al., \textbf{J.\ Ortiz}, ``VCHAR: Variance-Driven Complex Human Activity Recognition Framework with Generative Representation,'' \textit{arXiv} 2407.03291, 2024.

\item Y.\ Sun, J.\ Contreras, \textbf{J.\ Ortiz}, ``DFGauss: Dynamic Focused Masking for Autoregressive 3D Occupancy Prediction,'' \textit{NeurIPS}, 2025.

\item Z.\ Hussain et al., ``Non-Invasive Techniques for Monitoring Different Aspects of Sleep,'' \textit{ACM Trans.\ Comput.\ Healthcare}, 2022.

\item \textbf{J.\ Ortiz} et al., ``DeviceMien: Network Device Behavior Modeling for Identifying Unknown IoT Devices,'' \textit{ACM IoTDI}, 2019. \textbf{[Best Paper Finalist]}
\end{itemize}}

\section{Mentoring \& Service}

{\small\begin{itemize}
\item \textbf{Ph.D.\ graduates}: T.\ Chowdhury (2022), M.\ Aldeer (2023), T.\ Wu (2023) -- human activity sensing, multimodal learning, IoT
\item \textbf{Conference leadership}: Steering Cmte.\ Chair, ACM BuildSys (2024--25); General Chair, BuildSys (2022); TPC at ICLR, NeurIPS, MobiHoc, SenSys
\item \textbf{Broadening participation}: LatinX in AI Mentor; AAAI Undergrad.\ Consortium Coordinator; N2Women Co-Chair
\end{itemize}}

\end{document}
